\section{Introduction}


Text-to-image generative models have 
shown strong performance in recent years: 
models like Stable Diffusion~\citep{stablediffusion} and Dall-E~\citep{dalle} are capable of generating high quality and diverse images from arbitrary text prompts. However, a significant challenge faced by these models is that they rely \emph{solely} on text prompts alone for content control over the generation process, which is inadequate for many applications.  
Specifically, one of the most intuitive and user-friendly ways to exert control over the generation is to provide \emph{localization information}, which guides the models on where to generate specific elements within the image. Unfortunately, current pretrained models face limitations in their capability to perform localized generation. These limitations arise not only from their inability to incorporate location information as input but also from the inherent difficulties associated with compositionality, which is a known challenge for many multimodal foundation models~\citep{winoground}.




Existing methods addressing this issue typically fall into three main categories: training entirely new models~\citep{spade,pix2pix}, fine-tuning existing models with additional components such as task-specific encoders~\citep{gligen}, or strategically combining multiple samples into one~\citep{composablediffusion, multidiffusion}. All of these approaches often demand a substantial amount of training data, resources, and/or extended inference time, rendering them impractical for real-life applications due to their time and resource-intensive nature.
On the other hand, in a separate but related line of work, \citet{prompt2prompt} proposed Prompt-to-Prompt Image Editing, which edits generated images based on modified text prompts by manipulating cross attention maps in text-to-image generative models. Notably, this work also shows that cross attention layers play a pivotal role in controlling the spatial layout of generated objects associated with specific phrases in the prompts.


\begin{figure}[t]
  \centering
  \includegraphics[width=0.95\textwidth]{img/intro_figure.pdf}
  \caption[]{CAC as a plugin to existing methods for localized text-to-image generation. CAC improves upon diverse types of localization (bounding boxes, semantic segmentation maps and localized styles) with different base models (Stable Diffusion and GLIGEN). \footnotemark{}
  }
  \label{fig:examples} 
\end{figure}

In this work, we propose to use cross attention control (CAC) to provide pretrained text-to-image models with %
better open-vocabulary localization abilities.
As illustrated in Figure~\ref{fig:examples}, given a caption and localization information, such as bounding boxes and semantic segmentation maps, along with their corresponding text descriptions, we first construct a new text input by concatenating the caption and all prompts associated with the location information. We then compute the cross attention maps from this new text prompt and apply localization constraints to the cross attention maps according to the localization information. 
Our method does not require any additional training or model architecture modification like designing task-specific encoders. It also does not impose any language restrictions such as using a fixed set of vocabulary or a language parser. Moreover, it is highly portable and can be easily integrated into a single forward pass in any cross attention based text-to-image generation framework with only a few lines of code, thus demanding no extra inference time.




We develope a standardized suite of evaluation metrics for localized text-to-image generation tasks using off-the-shelf large pretrained recognition models~\citep{yolo,ssa,sam,glip}.
We apply CAC to various state-of-the-art baseline text-to-image generative models and experiment with different forms of localization information including bounding boxes and semantic segmentation maps. We demonstrate that CAC endows pretrained standard text-to-image models with new localized generation abilities, and furthermore, improves upon models specifically trained for localized generation. In addition, we show that with simple heuristics that spatially separate the components within text prompts, our method can significantly improve the compositional ability of text-to-image generative models.


\footnotetext{All shades of pink in the middle right example correspond to the prompt ``unicorn''.}